% Rozdział 6 — Podsumowanie i wnioski (budżet: 1,5 str.)
% 17.11.2026 (przegląd spójności): wersja robocza złożona wyłącznie z wyników
% i ograniczeń opisanych już w rozdz. 3–5; do akceptacji autora.
\chapter{Podsumowanie i~wnioski}
\label{ch:podsumowanie}

W~pracy zaprojektowano, wykonano i~zbadano prototyp miniaturowego
analizatora EDS w~technologii MEMS. Struktura składa się z~emitera
krzemowego z~ostrzem z~kompozytu CNT-PVP, bramki i~targetu
z~otworami, rozdzielonych szkłem borokrzemowym o~grubości
\SI{1}{\milli\metre} i~połączonych bondingiem anodowym w~atmosferze,
oraz z~magnesu ogniskującego NdFeB pod emiterem (\rozdzref{ch:projekt}).
Zestawiono stanowisko z~komorą próżniową, dwoma obwodami wysokiego
napięcia i~spektrometrem rentgenowskim (\rozdzref{ch:stanowisko}).

Najważniejsze wyniki badań (\rozdzref{ch:badania}):
\begin{itemize}
  \item emiter CNT wytwarza wiązkę elektronów przy napięciu ekstrakcyjnym
        powyżej ok. \SI{800}{\volt}; uzyskano prąd do \SI{22}{\micro\ampere}
        w~sesji bez magnesu i~do \SI{200}{\micro\ampere} w~sesji z~magnesem,
  \item z~magnesem wiązka tworzy na ekranie zwartą plamkę zamiast oświetlać
        całe okno, co potwierdza ogniskujące działanie pola magnetycznego;
        wpływu magnesu na prąd nie dało się oddzielić od stanu emitera,
        ponieważ odgazowano go tylko przed jedną sesją,
  \item wiązka wzbudza w~próbce promieniowanie charakterystyczne: przy
        $\UT = \SI{4}{\kilo\volt}$ zarejestrowano linie Mg, Al i~Si~\Ka,
        a~przy \SI{10}{\kilo\volt} zidentyfikowano pierwiastki czterech
        próbek (m.in. Cu, Pb, Au, Cr, Mn, Fe) po liniach K, L i~M,
  \item częstość zliczeń w~liniach charakterystycznych zmieniała się
        w~czasie pomiaru o~kilka do kilkunastu procent.
\end{itemize}

Cele pracy zrealizowano w~zakresie jakościowym. Ograniczeniami
prototypu są: brak hermetyczności, nieznane parametry magnesu, brak
skali na zdjęciach ekranu (średnicy plamki nie wyznaczono), brak
rejestracji prądu emisji podczas pomiarów widm oraz czas martwy
przekraczający \SI{50}{\percent} w~części pomiarów, który powoduje
powstawanie pików sumacyjnych. Linia tlenu nie była rejestrowana,
ponieważ zakres poniżej ok.~\SI{0.65}{\kilo\electronvolt} leży poza
czułością toru pomiarowego.

Wyniki wskazują, że struktura MEMS z~zimną katodą CNT może być źródłem
elektronów w~kompaktowym analizatorze EDS. Do zastosowania w~misji
kosmicznej potrzebne są: hermetyzacja przez bonding anodowy w~próżni
z~mikropompą albo praca w~naturalnej próżni ciała bez atmosfery,
napięcie przyspieszające \SIrange{10}{18}{\kilo\volt} do wzbudzania
linii K pierwiastków skałotwórczych łącznie z~żelazem, detektor
z~cienkim oknem do rejestracji tlenu oraz stabilizacja prądu wiązki
w~pętli sprzężenia zwrotnego.

W~dalszych pracach należy wykonać bonding w~próżni, zmierzyć indukcję
magnesu i~średnicę plamki ze wzorcem skali, rejestrować prąd emisji
podczas pomiarów widm, ograniczyć prąd wiązki tak, aby czas martwy nie
przekraczał ok.~\SI{30}{\percent}, oraz przeprowadzić analizę ilościową
z~użyciem wzorców.
